\chapter{Denoising fundamentals}

\section{A general solution to non-invertible problems}
%{{{

Non-invertible denoising problems can be approached by minimizing a
cost function
\begin{equation}
\mathbf{Y}^*=\underset{\mathbf{Y}}{\operatorname{arg\,min}} \, E(\mathbf{X},\mathbf{y}),
\end{equation}
where
\begin{equation}
E = D(\mathbf{X},\mathbf{Y}) + \lambda R(\mathbf{Y}),
\end{equation}
combines a data term $D$ and a regularizer term $R$.

The regularizer favors solutions based on prior assumptions about
uncorrupted images, and therefore, it helps to pick one restored image
$\mathbf{Y}^*$ from the set of uncorrupted images that might have
given rise to $\mathbf{X}$. $R$ traditionally encodes rather simple,
hand-crafted assumptions, often as simple as favoring smooth images by
penalizing strong differences of neighboring pixels
\cite{buchholz2019content}.

The data term $D$ encourages $\mathbf{Y}^*$ to be in-line with the
observation and our understanding of the physical nature of
distortions. Dependent on the severity of the distortion,
regularization has to be tuned by weighing it with a parameter
$\lambda$ \cite{buchholz2019content}.

%}}}

\section{Objectives}
%{{{

In microscopy, an optimal denoiser should effectively attenuate noise
while concurrently preserving crucial underlying biological
structures. Image denoising tries to restore a latent clean image from
an observed image degraded by noise \cite{BUCHHOLZ2019277}. Image
noise primarily arises during data acquisition, stemming from sensor
electronics, photon statistics, and variations in imaging conditions
such as illumination and scene dynamics. The resulting degradations
can be stochastic/uncorrelated (e.g., Poisson or Gaussian) or
structured/correlated (e.g., speckle or banding)
\cite{wang2025median2median}, both of which reduce the signal-to-noise
ratio, impair visual quality, and hinder downstream
analysis. Denoising methods are therefore needed to recover the
underlying clean image, either by modeling the corruption process or
by exploiting image priors such as self-similarity and sparsity.

An added challenge to this situation is the limited availability of
only a single noisy realization of these images, which complicates
both the evaluation of denoising algorithms (due to the absence of a
ground truth for comparative assessment) and the selection of
appropriate denoising parameters. 

%}}}

\section{Denoisers classification}
%{{{

Denoisers can be losely classified into:
\begin{enumerate}
\item \textbf{Algorithm-driven methods}: when their applicability
  depends fundamentaly on a list of computational steps
  ``hand-crafted'' optimized to improve the performance. Examples are
  signal-domain (such as Gaussian filters-based) denoisers,
  frequency-domain filtering (such as Wiener and wavelet-thresholding)
  denoisers, Variational (optimization-based) methods (such as \gls{TV}), MAP estimation, and Sparse coding -- confirmar
  esto), and signal-averaing methods (such as \gls{NLM} and
  \gls{BM3D}).
\item \textbf{Data-driven methods}: when the denoiser, although there
  is a list of computational steps, must \emph{learn} from the noisy
  data to improve the performance. Examples are all \gls{DL}
  techniques, that basically can be divided in \gls{AE} and
  transformers.
\end{enumerate}

\begin{enumerate}
  
\item Planar imaging-based denoising methods, that only a single 2D image to be applied.

  % Z. Y. M. Q. e. a. Qiao, C., “Zero-shot learning enables instant denoising
  % and super-resolution in optical fluorescence microscopy,” Nature Com-
  % munications, vol. 15, no. 4180, 2024.

  % P. R. S. A. e. a. Lequyer, J., “A fast blind zero-shot denoiser,” Nature
  % Machine Intelligence, vol. 4, p. 953–963, 2022.

\item Time-lapse imaging-based denoising methods, which leverage the
  temporal information present across multiple frames to distinguish
  between actual signal and random noise.
  
  % X. Li, G. Zhang, J. Wu, Y. Zhang, Z. Zhao, X. Lin, H. Qiao, H. Xie,
  % H. Wang, L. Fang et al., “Reinforcing neuron extraction and spike
  % inference in calcium imaging using deep self-supervised denoising,”
  % Nature Methods, pp. 1–6, 2021.

  % . H. S. N. O. P. L. . C. K. J´erˆome Lecoq, Michael Oliver, “Removing in-
  % dependent noise in systems neuroscience data using deepinterpolation,”
  % Nature Methods, vol. 18, pp. 1401–1408, 2021.

\end{enumerate}

%}}}
